\section*{ÍNDICE DE NOTACIONES}

\begin{flushleft}
$\mathbf{Tor}_n^A(P_\bullet,Q_\bullet)$: 6.3.1.

\medskip
$\mathscr{T}\!or_n^{\mathscr O_S}(\mathscr P_\bullet,\mathscr Q_\bullet)$,
$\mathscr{T}\!or_n^S(\mathscr P_\bullet,\mathscr Q_\bullet)$: 6.4.1 y 6.5.1.

\medskip
$\mathscr{T}\!or_n^S(\mathscr F,\mathscr G)$: 6.5.1.

\medskip
$\mathscr{T}\!or_q^S(\mathscr P_\bullet^{(1)},\mathscr P_\bullet^{(2)},\ldots,
\mathscr P_\bullet^{(m)})$: 6.5.15.

\medskip
$\mathbf H^{-n}(\mathfrak U^{(i)},\mathscr P_\bullet^{(i)})$,
$\mathbf H^{-n}(X^{(i)},\mathscr P_\bullet^{(i)})$: 6.6.1.

\medskip
$\mathbf{Tor}_n^A(L_{\bullet\bullet}^{(1)},L_{\bullet\bullet}^{(2)})$,
$\mathbf{Tor}_n^S(\mathfrak U^{(1)},\mathfrak U^{(2)};
\mathscr P_\bullet^{(1)},\mathscr P_\bullet^{(2)})$: 6.6.2.

\medskip
$\mathscr{T}\!or_n^S(\mathfrak U^{(1)},\mathfrak U^{(2)};
\mathscr F^{(1)},\mathscr F^{(2)})$,
$\mathscr{T}\!or_n^S(\mathfrak U^{(1)},\mathfrak U^{(2)};
\mathscr P_\bullet^{(1)},\mathscr P_\bullet^{(2)})$: 6.6.2.

\medskip
$\mathscr{T}\!or_n^S(f_1,f_2;\mathscr P_\bullet^{(1)},
\mathscr P_\bullet^{(2)})$: 6.6.8, 6.7.1 y 6.7.3.

\medskip
$\mathscr{T}\!or_n^S((f_i)_{i\in I};(\mathscr P_\bullet^{(i)})_{i\in I})$,
$\mathscr{T}\!or_n^S(f_1,\ldots,f_m;\mathscr P_\bullet^{(1)},\ldots,
\mathscr P_\bullet^{(m)})$: 6.7.11.

\medskip
$\mathbf{Ab}$, $\mathbf{Ab}_A$: 7.1.1.

\medskip
$T_{(B)}$, $T^{(B)}$, $T\otimes_A B$
($T$ funtor de $\mathbf{Ab}_A$ en $\mathbf{Ab}$): 7.1.3.

\medskip
$t_M$: 7.2.2.

\medskip
$T_{\mathfrak p}$: 7.1.4.

\medskip
$\mathscr T_\bullet^{Y'}(\mathscr M')$,
$\mathscr T_\bullet^U(\mathscr M|U)$: 7.7.2.

\medskip
$\operatorname{EP}(f,\mathscr P_\bullet;y)$,
$\operatorname{EP}(Y,\mathscr P_\bullet;y)$,
$\operatorname{EP}(\mathscr P_\bullet;y)$: 7.9.5.

\medskip
$\operatorname{PH}(f,\mathscr P_\bullet;y)$,
$\operatorname{PH}(\mathscr P_\bullet;y)$: 7.9.12.
\end{flushleft}

\clearpage
\oldpage[III]{214}
\section*{ÍNDICE TERMINOLÓGICO}

\begin{flushleft}
Característica de Euler-Poincaré de un complejo de módulos: 7.9.5.

\medskip
Cohomológicamente plano en un punto $y\in Y$ en dimensión $p$,
cohomológicamente plano sobre $Y$ en dimensión $p$, cohomológicamente plano
sobre $Y$ en el punto $y$, cohomológicamente plano sobre $Y$ (complejo de
$\mathscr O_X$-módulos planos sobre $Y$): 7.8.1.

\medskip
Complejos homotópicamente equivalentes: 6.1.4.

\medskip
Espacio anillado de dimensión cohomológica $\leq n$: 6.5.5.

\medskip
Extensión de escalares en un funtor covariante aditivo: 7.1.3.

\medskip
Haz de anillos de dimensión cohomológica $\leq n$: 6.5.5.

\medskip
Fórmula de Künneth: 6.7.8.

\medskip
Homológicamente plano en un punto $y\in Y$ en dimensión $p$,
homológicamente plano sobre $Y$ en dimensión $p$, homológicamente plano sobre
$Y$ en el punto $y$, homológicamente plano sobre $Y$ (complejo de
$\mathscr O_X$-módulos planos sobre $Y$): 7.8.1.

\medskip
Equivalencia homotópica: 6.1.4.

\medskip
Hipertor de dos complejos de $A$-módulos: 6.3.1.

\medskip
Hipertor local de dos complejos de módulos: 6.4.1.

\medskip
Hipertor global de dos complejos de módulos: 6.6.2, 6.6.8 y 6.7.3.

\medskip
Polinomio de Hilbert relativo a un complejo de módulos: 7.9.12.

\medskip
Sucesiones espectrales de Künneth: 6.7.3.
\end{flushleft}

\clearpage
\oldpage[III]{215}
\section*{TABLA DE MATERIAS}

\begin{flushleft}
\textsc{Capítulo III.} --- \textbf{Estudio cohomológico de los haces coherentes
(continuación)} \dotfill 5

\medskip
\S~6. Funtores «Tor» locales y globales; fórmula de Künneth \dotfill 5

6.1. Introducción \dotfill 5

6.2. Hipercohomología de los complejos de módulos sobre un preesquema \dotfill 6

6.3. Hipertor de dos complejos de módulos \dotfill 9

6.4. Funtores hipertor locales de complejos de módulos cuasicoherentes: caso
de los esquemas afines \dotfill 14

6.5. Funtores hipertor locales de complejos de módulos cuasicoherentes: caso
general \dotfill 16

6.6. Funtores de hipertor global para complejos de módulos cuasicoherentes y
sucesiones espectrales de Künneth: caso de base afín \dotfill 21

6.7. Funtores de hipertor global de complejos de módulos cuasicoherentes y
sucesiones espectrales de Künneth: caso general \dotfill 25

6.8. Sucesiones espectrales de asociatividad de los hipertor globales \dotfill 32

6.9. Las sucesiones espectrales de cambio de base en los hipertor globales
\dotfill 34

6.10. Estructura local de ciertos funtores cohomológicos \dotfill 39

\medskip
\S~7. Estudio del cambio de base en los funtores homológicos covariantes de
módulos \dotfill 43

7.1. Funtores de $A$-módulos \dotfill 43

7.2. Caracterización del funtor producto tensorial \dotfill 44

7.3. Criterios de exactitud de los funtores homológicos de módulos \dotfill 48

7.4. Criterios de exactitud para los funtores
$\mathrm H_i(P\otimes_A M)$ \dotfill 53

7.5. Caso de los anillos locales noetherianos \dotfill 58

7.6. Descenso de las propiedades de exactitud. Teorema de semicontinuidad y
criterio de exactitud de Grauert \dotfill 60

7.7. Aplicación a los morfismos propios: I. La propiedad de cambio de base
\dotfill 65

7.8. Aplicación a los morfismos propios: II. Criterios de planitud
cohomológica \dotfill 72

7.9. Aplicación a los morfismos propios: III. Invariancia de la
característica de Euler-Poincaré y del polinomio de Hilbert \dotfill 76

\begin{flushright}
\emph{Recibido el 15 de diciembre de 1962.}
\end{flushright}
\end{flushleft}

% La página impresa 216 está en blanco.
\clearpage
\thispagestyle{empty}
\null
\clearpage

\oldpage[III]{217}
\section*{ERRATAS Y ADENDAS}
\begin{center}
(Lista 2)

\medskip
A) \emph{Errores tipográficos}
\end{center}

\noindent\textbf{$(0_{\mathrm I},\,3.2.1)$} En la línea 2 de la p. 27, sustituir
$X$ por $U$.

\noindent\textbf{$(0_{\mathrm I},\,3.2.6)$} En la primera línea contando desde
abajo de la p. 27 y en la línea 1 de la p. 28, sustituir (3 veces)
$\mathfrak S_\lambda$ por $\mathscr H_\lambda$.

\noindent\textbf{$(0_{\mathrm I},\,3.4.5)$} Añadir un paréntesis delante de
3.4.5).

\noindent\textbf{$(0_{\mathrm I},\,4.2.1)$} En la línea 10 de la p. 39,
sustituir $\psi_*(A)$ por $\psi_*(\mathscr A)$.

\noindent\textbf{$(0_{\mathrm I},\,5.1.3)$} En la línea 16 de la p. 45,
sustituir $\mathscr F|V$ por $\mathscr F|U$.

\noindent\textbf{$(0_{\mathrm I},\,5.3.9)$} En la línea 12 de la p. 48, antes
de «tal que», añadir: sobre un entorno abierto $U$ de $x$.

\noindent\textbf{$(0_{\mathrm I},\,7.6.9)$} En la tercera línea contando desde
abajo de la p. 73, sustituir $I_\lambda$ por $J_\lambda$.

\noindent\textbf{(I, 1.1.15)} En la línea 3 de la p. 83, sustituir
$\mathfrak a$ por $\mathfrak a\ne A$.

\noindent\textbf{(I, 1.4.1)} En la segunda línea contando desde abajo de la
p. 90, sustituir $\widetilde N$ por $N$.

\noindent\textbf{(I, 2.3.1)} En la tercera línea contando desde abajo de la
p. 101, sustituir (0, 4.1.6) por (0, 4.1.7).

\noindent\textbf{(I, 2.3.2)} En la línea 11 de la p. 101, sustituir $B$ por
$B_s$ y $C$ por $C_t$.

\noindent\textbf{(I, 2.5.5)} En la línea 19 de la p. 104, sustituir
$S$-morfismo por $S$-preesquema.

\noindent\textbf{(I, 3.3.7)} En la línea 17 de la p. 109, sustituir
$Y_{(S)}$ por $Y_{(S')}$.

\noindent\textbf{(I, 3.4.5)} En la línea 9 de la p. 113, sustituir
$\mathrm s$ por $s$.

\noindent\textbf{(I, 3.4.8)} En la línea 4 de la p. 114, sustituir $p(x)$ por
$f(x)$.

\noindent\textbf{(I, 3.5.1)} En la línea 3 de la p. 115, sustituir
$1_X\times g$ por $1_{X'}\times g$.

\noindent\textbf{(I, 3.7.2)} En la línea 8 de la p. 119, sustituir el primer
$X$ por $X'$.

\noindent\textbf{(I, 4.2.2)} En la línea 8 de la p. 123, en la flecha vertical
izquierda del diagrama, sustituir $\alpha_{\psi(y)}$ por $\varphi_{\psi(y)}$.

\noindent\textbf{(I, 4.4.3)} En la línea 16 de la p. 126, sustituir «isomorfa»
por «isomorfas».

\noindent\textbf{(I, 4.5.5)} En la línea 17 contando desde abajo de la p. 127,
sustituir (4.2.4) por (4.2.5).

\noindent\textbf{(I, 5.1.4)} En la línea 14 contando desde abajo de la p. 128,
sustituir (2.1.7) por (2.1.8).

\noindent\textbf{(I, 5.3.13)} En la línea 20 de la p. 134, sustituir (4.2.4)
por (4.2.5).

\noindent\textbf{(I, 6.4.2)} En la primera línea contando desde abajo de la
p. 147, sustituir recubrimiento por recubrimiento finito.

\noindent\textbf{(I, 9.1.13)} En la línea 15 de la p. 171, sustituir
$p^{-1}(\mathscr F)$ por $p^*(\mathscr F)$.

\noindent\textbf{(I, 9.5.11)} En la línea 7 de la p. 179, sustituir $Y$ por
$Y'$.

\noindent\textbf{(I, 9.6.5)} En la línea 17 contando desde abajo de la p. 180,
sustituir $(0,\,4.1.4)$ por $(0,\,4.1.3)$.

\noindent\textbf{(I, 10.12.2)} En la línea 14 contando desde abajo de la p.
206, sustituir $(X_n)_{(S_n)}$ por $(X_n)_{(S_m)}$.

\clearpage
\oldpage[III]{218}
\noindent\textbf{(I), Índice terminológico:} En la línea 6 contando desde
abajo de la p. 219, sustituir 0, 4.1.4 por 0, 4.1.5. En las líneas 16, 17 y 18
contando desde abajo de la p. 222, sustituir I, 2.1.7 por I, 2.1.8; en la línea
19 contando desde abajo de la p. 222, sustituir I, 2.1.8 por I, 2.1.7.

\noindent\textbf{(II, 1.5.2)} En la línea 17 de la p. 12, insertar una $f$
junto a la flecha vertical izquierda del diagrama.

\noindent\textbf{(II, 4.2.7)} En la línea 16 contando desde abajo de la p. 75,
sustituir (III, 2.1.14) por (III, 2.1.13).

\noindent\textbf{(II, 5.2.1)} En la línea 16 contando desde abajo de la p. 97,
sustituir $X-U$ por $U$. En la línea 8 contando desde abajo de la p. 97,
sustituir $X'$ por $X_f$.

\noindent\textbf{(II, 5.3.6)} En la línea 16 de la p. 100, sustituir (4.6.18)
por (4.6.17).

\noindent\textbf{(II, 6.2.7)} En la línea 9 contando desde abajo de la p. 116,
sustituir cap. V por cap. IV.

\noindent\textbf{(II, 6.6.5)} En la línea 15 contando desde abajo de la p.
132, sustituir cap. V por cap. IV.

\noindent\textbf{(II, 7.4.12)} En la línea 1 de la p. 151, sustituir cap. V
por cap. III.

\noindent\textbf{(II, 8.1.4)} En la línea 15 de la p. 154, sustituir (III,
2.3.8) por (III, 2.3.7).

\noindent\textbf{(II, 8.10.5)} En la línea 17 contando desde abajo de la p.
187, sustituir $g_{(j)x}$ por $g_{j(x)}$.

\noindent\textbf{$(0_{\mathrm{III}},\,10.2.7)$} En la línea 1 de la p. 20,
sustituir $u$-ádico por $\mathfrak m$-ádico. En la línea 13 contando desde abajo
de la p. 20, sustituir $k$ por $k_i$.

\noindent\textbf{$(0_{\mathrm{III}},\,11.1.1)$} En la línea 18 de la p. 23,
sustituir $\operatorname{Im}(d_r^{p+r,q-r+1})$ por
$\operatorname{Im}(d_r^{p-r,q+r-1})$.

\noindent\textbf{$(0_{\mathrm{III}},\,11.8.6)$} En la línea 14 contando desde
abajo de la p. 45, sustituir las flechas $\to$ por $\rightsquigarrow$ (4 veces).

\noindent\textbf{$(0_{\mathrm{III}},\,12.3.4)$} En la línea 4 de la p. 61,
sustituir $\rightsquigarrow$ por $\to$.

\noindent\textbf{$(0_{\mathrm{III}},\,13.2.4)$} En la segunda línea contando
desde abajo de la p. 67, sustituir $\varinjlim$ por $\varprojlim$ (2 veces).

\noindent\textbf{(III, 1.1.7)} En la línea 15 de la p. 85, sustituir $A^r$ por
$\bigwedge(A^r)$ (2 veces).

\noindent\textbf{(III, 1.2.4)} En la línea 7 de la p. 87, sustituir
$\Gamma(U,\mathscr F)$ por $\Gamma(U,\mathscr O_X)$ y sustituir $\mathscr G$
por $\mathscr F$. En la línea 8 de la p. 87, sustituir
$\mathrm H^p(\mathfrak U,\mathscr G)$ por $\mathrm H^p(\mathfrak U,\mathscr F)$.

\noindent\textbf{(III, 2.5.3.1)} En la línea 3 de la p. 111, sustituir $n>0$
por $n\geq0$. En la línea 8 de la p. 111, sustituir $-r\leq n\leq0$ por
$-r\leq n<0$.

\noindent\textbf{(III, 3.1.2)} En la línea 13 de la p. 116, sustituir
$\mathscr G$ por $\mathscr G\in\mathbf K'$.

\noindent\textbf{(III, 4.3.6)} En la línea 7 de la p. 133, sustituir
$K'=K\otimes_A\widehat A$ por $K'\supset K\otimes_A\widehat A$. En la línea 16
de la p. 133, sustituir $R(X)$ por $R(Y)$.

\noindent\textbf{(III, 4.4.12)} En la línea 10 de la p. 138, sustituir cap. V
por cap. IV.

\noindent\textbf{(III, 4.6.6)} En la línea 19 de la p. 142, sustituir «cap. V»
por «en un parágrafo posterior».

\begin{center}
B) \emph{Modificaciones del texto}\footnote{Para facilitar la búsqueda de las
referencias, las modificaciones pertenecientes a la lista de erratas insertada
en el capítulo N serán designadas en lo sucesivo por $\mathbf{Err}_{\mathrm N}$
seguido de un número.}
\end{center}

\noindent\textbf{$(\mathbf{Err}_{\mathrm{III}},\,1)$} En
$(0_{\mathrm I},\,5.1.3)$, en la línea 18 de la p. 45, sustituir «suma directa»
por «suma directa finita».

\noindent\textbf{$(\mathbf{Err}_{\mathrm{III}},\,2)$} En
$(0_{\mathrm I},\,5.4.3)$, en las líneas 8 a 4 contando desde abajo de la p. 49,
sustituir desde «En el caso general\ldots{}» por el texto siguiente:

En el caso general en que $X$ es un espacio anillado tal que $\mathscr O_x$ es
un anillo \emph{local} para
\oldpage[III]{219}
todo $x\in X$, si $\mathscr L$ es un $\mathscr O_X$-módulo \emph{de tipo finito}
tal que existe un $\mathscr O_X$-módulo $\mathscr F$ para el cual
$\mathscr L\otimes_{\mathscr O_X}\mathscr F$ es \emph{isomorfo a}
$\mathscr O_X$, el razonamiento anterior muestra que, para todo $x\in X$,
$\mathscr L_x$ es isomorfo a $\mathscr O_x$. De ello se deduce que $\mathscr L$
es entonces \emph{invertible}. En efecto, para todo $x\in X$, sea $U$ un
entorno abierto de $x$ tal que $\mathscr L|U$ esté engendrado por $n$ secciones
$s_i\ (1\leq i\leq n)$ sobre $U$ (5.2.1). Se puede suponer, por ejemplo, que
$s=s_1$ es tal que $s_x\ne0$ y, como $\mathscr L_x$ es isomorfo a
$\mathscr O_x$, existe para cada $i$ una sección $t_i$ de $\mathscr O_X$ sobre
un entorno abierto $V_i\subset U$ de $x$ tal que
$(t_i)_x s_x=(s_i)_x$. Por consiguiente, existe un entorno abierto
$V\subset\bigcap_{i=1}^nV_i$ de $x$ tal que $t_i s=s_i$ en $V$; dicho de otro
modo, $\mathscr L|V$ está engendrado por \emph{una sola} sección $s|V$. Además,
si $z$ es una sección, sobre un abierto $W\subset V$, del núcleo del
homomorfismo $\mathscr O_X|V\to\mathscr L|V$ definido por $s$ (5.1.1), $z_y$
anula $\mathscr L_y$ para todo $y\in W$, luego $z_y=0$ por hipótesis y, por
consiguiente, $z=0$, lo que termina de probar nuestra afirmación. Además, la
consideración del producto tensorial
$\mathscr L^{-1}\otimes\mathscr L\otimes\mathscr F$ muestra de inmediato que
$\mathscr F$ es isomorfo a $\mathscr L^{-1}$.

\noindent\textbf{$(\mathbf{Err}_{\mathrm{III}},\,3)$} En
$(0_{\mathrm I},\,7.2.4)$, es preciso imponer la condición de que las potencias
$\mathfrak J^n$ sean ideales \emph{cerrados} en el anillo admisible $A$ para
que la proposición sea correcta. La demostración dada es incorrecta, y el
enunciado rectificado resulta de Bourbaki, \emph{Top. gén.}, cap. III,
3.\textsuperscript{a} ed., \S~3, n.\textsuperscript{o}~5, cor. 1 de la prop. 9.
Asimismo, es preciso suponer que los $\mathfrak J^n$ son \emph{cerrados} en $A$
en los enunciados $(0_{\mathrm I},\,7.2.5)$ y $(0_{\mathrm I},\,7.2.6)$.

\noindent\textbf{$(\mathbf{Err}_{\mathrm{III}},\,4)$} Después de la
proposición (I, 2.2.5), añadir: con las notaciones de (I, 2.2.5), si
$\mathfrak J$ es un ideal de $A$, se denota por
$\widetilde{\mathfrak J\mathscr F}$ el sub-$\mathscr O_X$-módulo
$\mathfrak J\widetilde{\mathscr F}$ definido en $(0,\,4.3.5)$.

\noindent\textbf{$(\mathbf{Err}_{\mathrm{III}},\,5)$} En (I, 3.4.5), en la
primera línea contando desde abajo de la p. 112, sustituir «\emph{un cuerpo}
$K$» por «\emph{un cuerpo algebraicamente cerrado} $K$». En la línea 1 de la
p. 113, sustituir «\emph{extensión}» por «\emph{extensión algebraicamente
cerrada}». En las líneas 1 a 4 de la p. 113, sustituir desde «$K$ será
llamado\ldots{}» por el texto siguiente:

Para todo punto de $X$ con valores en un cuerpo $K$, $K$ se llamará el
\emph{cuerpo de valores} del punto correspondiente y, si $x$ es la localidad
de este punto, se dice también que este último está \emph{localizado en} $x$.
Se define así una aplicación $X(K)\to X$ que hace corresponder a un punto con
valores en $K$ su localidad.

En las líneas 7 y 16 de la p. 113, suprimir «geométrico»; en las líneas 10 y 13
de la p. 113, suprimir «geométricos». En las líneas 10 y 11 contando desde
abajo de la p. 115, suprimir «\emph{geométrico}».

\noindent\textbf{$(\mathbf{Err}_{\mathrm{III}},\,6)$} Al final de la
demostración de (I, 3.6.1), en la línea 12 contando desde abajo de la p. 117,
añadir: Si se pone
$X'=X\times_Y\operatorname{Spec}(\mathscr O_y/\mathfrak a_y)$, entonces, para
todo punto $x\in X'$, identificado por $p$ con un punto de $X$, se tiene
$\mathscr O_{X',x}=\mathscr O_{X,x}/\mathfrak a_y\mathscr O_{X,x}$. Puesto que
la cuestión es local en $X$ y $Y$, se puede suponer que
$X=\operatorname{Spec}(B)$, $Y=\operatorname{Spec}(A)$, y se tiene
$(B/\mathfrak a_yB)_x=B_x/\mathfrak a_yB_x$ por planitud $(0,\,1.3.2)$.

\noindent\textbf{$(\mathbf{Err}_{\mathrm{III}},\,7)$} En (I, 5.1.1), en la
tercera línea contando desde abajo de la p. 127, sustituir
«$\mathscr O_X$-módulo cuasicoherente» por «\emph{ideal cuasicoherente de}
$\mathscr B$».

\noindent\textbf{$(\mathbf{Err}_{\mathrm{III}},\,8)$} Después de (I, 5.1.10),
en la línea 16 de la p. 131, añadir:

\emph{Observación (5.1.11).} --- Una ligera adaptación del razonamiento hecho
en (5.1.9) muestra que la conclusión sigue siendo válida \emph{sin suponer a
priori que $X$ sea un preesquema}, sino

\oldpage[III]{220}
suponiendo solamente que $(X,\mathscr O_X)$ sea un espacio anillado en anillos
locales, que $\mathscr J$ sea un ideal de $\mathscr O_X$ tal que
$\mathscr J^n=0$, que el espacio anillado
$X_0=(X,\mathscr O_X/\mathscr J)$ sea un esquema afín y, por último, que los
ideales $\mathscr J^k/\mathscr J^{k+1}$ sean $\mathscr O_{X_0}$-módulos
cuasicoherentes. Basta, en efecto, utilizar (1.8.1) (cf.
$(\mathbf{Err}_{\mathrm{II}})$) en lugar de (2.2.4) en el razonamiento.

\noindent\textbf{$(\mathbf{Err}_{\mathrm{III}},\,9)$} En (I, 5.3.1), en la
línea 11 de la p. 132, insertar, después de «o $\Delta_X$», «o $\Delta_\varphi$
si $\varphi:X\to S$ es el morfismo estructural».

\noindent\textbf{$(\mathbf{Err}_{\mathrm{III}},\,10)$} En (I, 5.3.9), la
demostración es insuficiente, pues no prueba que $\Delta_X(X)$ sea localmente
cerrado en $X\times_SX$. Para obtener una demostración correcta, basta utilizar
(4.2.4, $a$)): para todo $x\in X$ y todo entorno afín $U$ de $x$ en $X$,
$U\times_SU$ es un entorno afín de $\Delta_X(x)$; teniendo en cuenta (5.3.16)
(cuya demostración solo utiliza la definición (5.3.1.1)), se reduce a demostrar
(5.3.9) cuando $S=\operatorname{Spec}(B)$ y $X=\operatorname{Spec}(A)$ son
esquemas afines; entonces resulta claro, en virtud de (5.3.1.1), que $\Delta_X$
corresponde al homomorfismo canónico $A\otimes_BA\to A$ que transforma
$x\otimes y$ en $xy$; como este homomorfismo es sobreyectivo, $\Delta_X$ es en
este caso una inmersión cerrada (4.2.3), lo que termina la demostración.

\noindent\textbf{$(\mathbf{Err}_{\mathrm{III}},\,11)$} En (I, 7.1.15) y
(I, 7.1.16), en las líneas 7 y 15 contando desde abajo de la p. 158, suprimir
«\emph{geométricos}».

\noindent\textbf{$(\mathbf{Err}_{\mathrm{III}},\,12)$} Sustituir (I, 7.4.7)
por: Se extienden (por abuso de lenguaje) las definiciones de (7.4.1) al caso
en que $X$ es un preesquema \emph{reducido} tal que todo punto admite un
entorno abierto que solo tiene un número \emph{finito} de componentes
irreducibles; se deduce entonces de (7.3.4) y (7.4.6) que, para un
$\mathscr O_X$-módulo cuasicoherente de tipo finito $\mathscr F$, decir que
$\mathscr F$ es un \emph{haz de torsión} equivale a decir que
$\operatorname{Supp}(\mathscr F)$ \emph{no contiene ninguna componente
irreducible de $X$}.

\noindent\textbf{$(\mathbf{Err}_{\mathrm{III}},\,13)$} En (I, 10.11.7), en
las líneas 17 a 23 de la p. 205, el final de la demostración desde «Queda por
probar\ldots{}» es innecesario en virtud de (10.11.6), pues los haces
considerados son coherentes por $(0,\,5.3.5)$.

\noindent\textbf{$(\mathbf{Err}_{\mathrm{III}},\,14)$} En (II, 1.7.8),
después de la línea 10 contando desde abajo de la p. 16, añadir: cuando
$\mathscr E=\mathscr O_S^n$, se escribe también $\mathbf V_S^n$ en lugar de
$\mathbf V(\mathscr E)$; si además $S=\operatorname{Spec}(A)$ es afín, se
escribe $\mathbf V_A^n$ en lugar de $\mathbf V_S^n$; se tiene
$\mathbf V_A^n=\operatorname{Spec}(A[T_1,\ldots,T_n])$, donde los $T_i$ son
indeterminadas.

\noindent\textbf{$(\mathbf{Err}_{\mathrm{III}},\,15)$} En (II, 1.7.10), en
la línea 10 de la p. 17, suprimir «\emph{geométricos}»; en las líneas 12 y
16--17 de la p. 17, suprimir «geométrico».

\noindent\textbf{$(\mathbf{Err}_{\mathrm{III}},\,16)$} En (II, 1.7.12),
añadir: Se pone asimismo:
\[
  \mathbf V(\mathscr O_S^n)=\mathbf V_S^n=S[T_1,\ldots,T_n].
\]

\noindent\textbf{$(\mathbf{Err}_{\mathrm{III}},\,17)$} En (II, 4.2.6), en la
línea 7 de la p. 75, suprimir «\emph{geométricos}»; en la línea 8 de la p. 75,
suprimir «\emph{geométrico}».

\noindent\textbf{$(\mathbf{Err}_{\mathrm{III}},\,18)$} En (II, 4.6.13), en la
línea 7 de la p. 92, es preciso especificar el razonamiento, pues, con las
notaciones de (4.4.10), hay que probar aquí que la inmersión $\Gamma_f$ es
\emph{cuasicompacta}, a fin de aplicar (i bis). En virtud de (4.6.4), para
probar que $\mathscr L$ es amplio para $f$, se puede limitar al caso
en que $Y$ es afín. Obsérvese, por otra parte, que en cada una de las hipótesis
de (v), $f$ es cuasicompacto (I, 6.6.4). Por otra parte, si $g$ es
\oldpage[III]{221}
separado, $\Gamma_f$ es una inmersión cerrada (I, 5.4.3), luego cuasicompacta
(I, 6.6.4). Si, por el contrario, $X$ es localmente noetheriano, como $Y$ es
afín y $f$ cuasicompacto, el espacio subyacente a $X$ es cuasicompacto, luego
noetheriano, y se puede aplicar de nuevo (I, 6.6.4) para probar que $\Gamma_f$
es cuasicompacta.

\noindent\textbf{$(\mathbf{Err}_{\mathrm{III}},\,19)$} El enunciado de (II,
5.2.2) es incorrecto, pues las condiciones $b)$, $c)$ y $c')$ no son locales
en $Y$ cuando no se sabe si, para un abierto $U$ de $Y$, un
$(\mathscr O_X|f^{-1}(U))$-módulo cuasicoherente es la restricción a
$f^{-1}(U)$ de un $\mathscr O_X$-módulo cuasicoherente. Por tanto, es preciso
sustituir en la línea 16 de la p. 98 «\emph{Sea $f:X\to Y$ un morfismo separado
y cuasicompacto}» por: «\emph{Sean $X$, $Y$ dos preesquemas tales que $X$ sea
un esquema o que el espacio subyacente a $X$ sea localmente noetheriano, y sea
$f:X\to Y$ un morfismo cuasicompacto.}» En la demostración, se observará que
las hipótesis implican que $f$ es \emph{separado} cuando se supone que $X$ es
un esquema, en virtud de (I, 5.5.5 y 5.5.8); por tanto, se puede aplicar
(I, 9.2.2) en ambos casos.

\noindent\textbf{$(\mathbf{Err}_{\mathrm{III}},\,20)$} En (II, 6.2.3),
añadir, después de la línea 18 contando desde abajo de la p. 115:

\emph{Se dice que un morfismo $f:X\to Y$ es cuasifinito en un punto $x\in X$
si existen un entorno abierto afín $V$ de $y=f(x)$ y un entorno abierto afín
$U$ de $x$ tales que $f(U)\subset V$ y que el morfismo $U\to V$, restricción de
$f$, sea cuasifinito. Se dice que un morfismo $f:X\to Y$ es localmente
cuasifinito si es cuasifinito en todo punto de $X$.}

\noindent\textbf{$(\mathbf{Err}_{\mathrm{III}},\,21)$} En (II, 6.4.3), la
demostración es incorrecta, pues los elementos de un sub-$A$-módulo de tipo
finito de $K$ no son necesariamente enteros sobre $A$. Sustituir las 9 últimas
líneas de la p. 121 por el texto siguiente: Como se puede suponer que $E$ es
sin torsión, luego fiel, $E$ es también un $A[u]$-módulo fiel ($A$ está
inmerso en el anillo de endomorfismos de $E$). Como $E$ es un $A$-módulo de
tipo finito, se deduce que $u$ es \emph{entero} sobre $A$ (Bourbaki,
\emph{Alg. comm.}, cap. V, \S~1, n.\textsuperscript{o}~1, lema 1). Se concluye
inmediatamente que los valores propios de $u\otimes1$ (en una clausura
algebraica de $K$) son elementos \emph{enteros} sobre $A$, y por tanto también
lo son los $\sigma_i(u)$.

\noindent\textbf{$(\mathbf{Err}_{\mathrm{III}},\,22)$} En
$(0_{\mathrm{III}},\,9.1.1)$, en la línea 18 contando desde abajo de la p. 12,
suprimir «abierta».

\noindent\textbf{$(\mathbf{Err}_{\mathrm{III}},\,23)$} En
$(0_{\mathrm{III}},\,11.7.3)$, en la línea 10 de la p. 43, después de
$\mathsf C''$, añadir: tal que, para todo objeto proyectivo $P$ de
$\mathsf C'$ (respectivamente, todo objeto proyectivo $P'$ de $\mathsf C''$),
el funtor $A'\rightsquigarrow T(P,A')$ (respectivamente,
$A\rightsquigarrow T(A,P')$) sea exacto en $\mathsf C'$ (respectivamente, en
$\mathsf C$).

\noindent\textbf{$(\mathbf{Err}_{\mathrm{III}},\,24)$} El enunciado de
$(0_{\mathrm{III}},\,13.7.7)$ es inexacto y debe modificarse como sigue: en la
línea 6 de la p. 78, suprimir «y
$(\mathrm R^{n+1}T(A_k))_{k\in\mathbf Z}$»; en la línea 7 de la p. 78, sustituir
«$\mathrm R'^nT(A)$ es un $S$-módulo de tipo finito» por
«$\mathrm R'^nT(A)$ y $\mathrm R'^{n+1}T(A)$ son $S$-módulos de tipo finito»;
en las líneas 12--13 de la p. 78, suprimir «y $p+q=n+1$». En la demostración,
en la línea 23 de la p. 78, suprimir «y para $n+1$».

\noindent\textbf{$(\mathbf{Err}_{\mathrm{III}},\,25)$} En (III, 1.4.15), en
la línea 6 contando desde abajo de la p. 92, sustituir «\emph{de tipo finito}»
por «\emph{cuasicompacto}».

\noindent\textbf{$(\mathbf{Err}_{\mathrm{III}},\,26)$} En (III, 2.2.4), en la
línea 1 de la p. 102, sustituir «\emph{que los soportes de $\mathscr F$ y de
$\mathscr H$ sean propios sobre $Y$}» por «\emph{que el soporte de
$\operatorname{Im}(\mathscr F\to\mathscr G)$ sea propio sobre $Y$}». En la
demostración, sustituir el texto de las líneas 10 a 16, desde «Dicho
esto\ldots{}», por:

\oldpage[III]{222}
Pongamos
$\mathscr G_1=\operatorname{Im}(\mathscr F\to\mathscr G)
=\operatorname{Ker}(\mathscr G\to\mathscr H)$, que es coherente. Como se tiene
la sucesión exacta
$0\to\mathscr G_1(n)\to\mathscr G(n)\to\mathscr H(n)$ y el funtor $f_*$ es
exacto por la izquierda, la sucesión
$0\to f_*(\mathscr G_1(n))\to f_*(\mathscr G(n))\to f_*(\mathscr H(n))$ es
exacta para todo $n$, y basta por tanto mostrar que, para $n$ suficientemente
grande, la sucesión $f_*(\mathscr F(n))\to f_*(\mathscr G_1(n))\to0$ es exacta.
Dicho de otro modo, se puede limitar al caso en que $\mathscr H=0$ y en que el
soporte de $\mathscr G$ es propio sobre $Y$, luego \emph{cerrado en $Z$}
(II, 5.4.10); por consiguiente, $\mathscr G'=i_*(\mathscr G)$ es un
$\mathscr O_Z$-módulo coherente tal que $\mathscr G'|X=\mathscr G$. Se sabe
(I, 9.4.3) que existe un $\mathscr O_Z$-módulo coherente $\mathscr F'$ tal que
$\mathscr F'|X=\mathscr F$. Además, como $X$ es abierto en $Z$ y
$\operatorname{Supp}(\mathscr G)\subset X$ es cerrado en $Z$, es inmediato que
se define un homomorfismo sobreyectivo $u':\mathscr F'\to\mathscr G'$ de haces
tal que $u'|X$ sea el homomorfismo sobreyectivo dado
$u:\mathscr F\to\mathscr G$, tomando $u'|U=0$ para todo abierto $U$ de $Z$ que
no interseque $\operatorname{Supp}(\mathscr G)$ y $u'|U=u|U$ para todo abierto
$U\subset X$. Dicho esto, se ve, como al comienzo de la demostración de
(2.2.2), que se puede limitar al caso en que $Y$ es afín, y se trata por tanto
de mostrar que, para $n$ suficientemente grande, el homomorfismo
$\Gamma(u):\Gamma(X,\mathscr F(n))\to\Gamma(X,\mathscr G(n))$ es sobreyectivo.
Ahora bien, se tiene el diagrama conmutativo
\[
\xymatrix@C=35mm@R=12mm{
\Gamma(Z,\mathscr F'(n)) \ar[r]^{\Gamma(u')} \ar[d]_v &
\Gamma(Z,\mathscr G'(n)) \ar[d]^w \\
\Gamma(X,\mathscr F(n)) \ar[r]_{\Gamma(u)} & \Gamma(X,\mathscr G(n))
}
\]
donde $v$ y $w$ son los homomorfismos de restricción. La definición de
$\mathscr G'$ muestra además que $w$ es \emph{biyectivo}; por otra parte, en
virtud de (2.2.3), $\Gamma(u')$ es \emph{sobreyectivo} para $n$ suficientemente
grande, luego lo mismo vale para $\Gamma(u)$.

\noindent\textbf{$(\mathbf{Err}_{\mathrm{III}},\,27)$} En (III, 2.2.5),
después de la línea 7 contando desde abajo de la p. 102, añadir:

(iii) Bajo las hipótesis de (2.2.4) relativas a $X$, $Y$, $f$ y $\mathscr L$,
es inmediato que, si $\mathscr H$ es un $\mathscr O_X$-módulo coherente cuyo
soporte es \emph{propio sobre $Y$}, un razonamiento análogo al de (2.2.4)
muestra que existe un entero $N$ tal que, para $n\geq N$, se tiene
$\mathrm R^if_*(\mathscr H(n))=0$ para todo $i>0$. Se deduce, por la sucesión
exacta de cohomología, que, si $u:\mathscr F\to\mathscr G$ es un homomorfismo
de $\mathscr O_X$-módulos coherentes tal que $\operatorname{Ker}(u)$ y
$\operatorname{Coker}(u)$ tienen soportes \emph{propios sobre $Y$}, entonces
existe $N$ tal que, para $n\geq N$, el homomorfismo correspondiente
$\mathrm R^if_*(\mathscr F(n))\to\mathrm R^if_*(\mathscr G(n))$ sea
\emph{biyectivo} para todo $i>0$.

\noindent\textbf{$(\mathbf{Err}_{\mathrm{III}},\,28)$} En (III, 4.4.9), en la
línea 17 contando desde abajo de la p. 137, sustituir «\emph{Sean $Y$ un
preesquema íntegro localmente noetheriano}» por «\emph{Sean $X$ e $Y$ dos
preesquemas íntegros localmente noetherianos}»; en la línea 16 contando desde
abajo de la p. 137, sustituir «\emph{de tipo finito}» por «\emph{localmente de
tipo finito}». La demostración permanece esencialmente inalterada, pues
$f^{-1}(y)$ es localmente de tipo finito sobre $k(y)$ y, por tanto, sigue
siendo discreto.

\noindent\textbf{$(\mathbf{Err}_{\mathrm{III}},\,29)$} En (I, 9.3.4), en la
línea 19 contando desde abajo de la p. 173, sustituir «\emph{noetheriano}» por
«\emph{cuasicompacto}»; en las líneas 18 y 19 contando desde abajo de la p.
173, sustituir «\emph{coherente}» por «\emph{cuasicoherente de tipo finito}».
En la demostración, se reduce al caso en que

\oldpage[III]{223}
$X=\operatorname{Spec}(A)$, $\mathscr F=\widetilde M$, donde $M$ es un
$A$-módulo de tipo finito, $\mathscr J=\widetilde{\mathfrak J}$, donde
$\mathfrak J$ es un ideal de tipo finito de $A$; el resto del razonamiento
permanece entonces inalterado.

\noindent\textbf{$(\mathbf{Err}_{\mathrm{III}},\,30)$} En (I, 9.3.5),
sustituir las líneas 1 a 7 contando desde abajo de la p. 173 por el texto
siguiente:

\emph{Proposición (9.3.5). --- Sean $X$ un preesquema y $\mathscr F$ un
$\mathscr O_X$-módulo cuasicoherente de tipo finito. Entonces existen un
subpreesquema cerrado $Y$ de $X$, cuyo espacio subyacente es igual a
$\operatorname{Supp}(\mathscr F)$, y un $\mathscr O_Y$-módulo cuasicoherente de
tipo finito $\mathscr G$ tales que, si $j:Y\to X$ es la inyección canónica,
$\mathscr F$ sea isomorfo a $j_*(\mathscr G)$.}

Bastará mostrar que el ideal $\mathscr J$ de $\mathscr O_X$,
\emph{anulador de $\mathscr F$}, es \emph{cuasicoherente}; se tomará entonces
como $Y$ el subpreesquema cerrado de $X$ definido por $\mathscr J$ (I, 4.1.2),
y, como $\mathscr J\mathscr F=0$, $\mathscr F$ es un
$(\mathscr O_X/\mathscr J)$-módulo y se resolverá la cuestión tomando
$\mathscr G=j^*(\mathscr F)$. Para ver que $\mathscr J$ es cuasicoherente, se
puede (puesto que la cuestión es local) limitar al caso en que
$X=\operatorname{Spec}(A)$, $\mathscr F=\widetilde M$, donde $M$ es un
$A$-módulo engendrado por un número finito de elementos
$x_i\ (1\leq i\leq r)$; el ideal $\mathscr J$ es entonces la intersección de
los anuladores de los $x_i$. Pero el anulador de $x_i$ es el núcleo del
homomorfismo $\mathscr O_X\to\mathscr F$ correspondiente al homomorfismo
$s\mapsto sx_i$ de $A$ en $M$; es, por tanto, un ideal cuasicoherente (I,
4.1.1), y toda intersección finita de tales ideales es asimismo un ideal
cuasicoherente (I, 1.3.10).
